% =====================================================================================================
% Mục mới (bản nháp 25/10/2026) theo kiểm toán 23/10 §6-P1, §6-P3, §7: so sánh với Elmi Ordinary trên dữ liệu Via
% trong bộ mô phỏng, và kiểm chứng bộ mô phỏng bằng RideBound thật. Được \input trong 08-chuong4.tex, trước mục Thảo luận.
% Số liệu: tables_2x3.py (tính lại độc lập từ wf14/results/main.ndjson và wf15/results/confirm.ndjson), P5-RESULT.md,
% BAO-CAO-GIAI-DOAN-1.md, P7-RESULT.md (ridebound-scratchpad, 25/10/2026).
% =====================================================================================================
\section{So sánh với Elmi Ordinary trên dữ liệu Via trong bộ mô phỏng}
\label{muc:c4-via}

Các mục trước chạy RideBound thật trên những kịch bản Manhattan nhỏ, nơi phần cộng dồn và phần hai chiều của luật C
hiếm khi có việc. Mục này đặt luật C cạnh đối chứng của Elmi và cộng sự~\cite{elmi2026} ở mật độ cao hơn nhiều, bằng một
bộ mô phỏng. Mọi số trong mục là số của bộ mô phỏng, không phải của RideBound; phần kiểm chứng bộ mô phỏng bằng RideBound
thật ở mục~\ref{muc:c4-cauno}.

\subsection{Dữ liệu và cấu hình}
\label{muc:c4-via-dulieu}

Nhu cầu lấy từ bản ghi chuyến của các hãng xe lớn mà Ủy ban Taxi và Limousine New York (TLC) công bố, phần của Via
(mã hãng HV0004), đón và trả trong Manhattan, năm 2019. Mỗi kịch bản là một khung hai giờ (8--10 giờ hoặc 17--19 giờ)
của một ngày thường, gồm 512 yêu cầu rút ngẫu nhiên có hạt giống cố định từ các yêu cầu Via của khung đó, 62 xe bốn chỗ,
thời gian đi tối đa gấp 2 lần đi thẳng tính ở đầu khung, thời gian dừng đón trả 0 giây. Điểm đón và trả là một nút ngẫu
nhiên của mạng đường Manhattan trong khu vực TLC tương ứng. Hệ số tắc đường theo khung 15 phút được ước lượng từ chuyến
taxi vàng, chung cho cả thành phố. Tập so sánh gồm 60 ngày thường của 10 tháng (6 ngày mỗi tháng), mỗi ngày hai khung và
ba lần rút mẫu, tổng 360 kịch bản.
% EV: sp/wf14/scripts/gen_cells_rs.js:5-9, 35, 96, 104-108; sp/wf15/scripts/shlib.js:17 (maxRidePermille 2000, dwellMs 0, nVeh 62); sp/wf15/work/days_confirm.csv (60 ngày)

Để chắc dữ liệu đầu vào là chuyến thật, hai tháng (3/2019 và 9/2019, 12 ngày, 118.075 yêu cầu) được trích lại từ tệp
gốc của TLC và so với đầu vào của bộ mô phỏng: không lệch dòng nào theo cả hai chiều, kể cả số thứ tự yêu cầu dùng để
gán nút. Dữ liệu này chỉ gồm chuyến đã hoàn thành; không có yêu cầu bị hủy hay bị từ chối.
% EV: p7-tlc-2026-10-25/P7-RESULT.md (EXCEPT ALL 0/0, rid theo request_datetime, pu, dz, trip_miles, trip_time)

\subsection{Đối chứng}
\label{muc:c4-via-doichung}

Cách siết đều của Elmi và cộng sự được hình thức hóa như sau: mỗi khách nhận ngân sách $B_r = \mathrm{round}(\alpha \cdot
\delta_r)$, với $\delta_r$ là độ dư giờ trả lúc hứa đầu tiên (giờ đón dự kiến cộng thời gian đi tối đa, trừ giờ trả đã
hứa) và $\alpha = 35\%$; giờ trả mới phải thỏa $p \le p_0 + B_r$. Ở bảng gốc của bài, phía trễ tương ứng khoảng 20\% độ
dư, nên cách hình thức hóa này lỏng hơn ở phía trễ. Bài Elmi giả định thời gian đi cố định, nên không nói hàng rào xử lý
thế nào khi tắc đường đẩy giờ trả dự kiến $e$ vượt nó. Khóa luận xét hai cách mở rộng: hàng rào chặt (EO35), giữ nguyên
tuyến của xe khi không còn phương án hợp lệ; và hàng rào nới (EOHZ35), $p \le \max(p_0 + B_r, e)$, chính là dạng hạn chót
nới ở mục~\ref{muc:doichung}. Hàng rào nới là đối chứng chính; hàng rào chặt luôn được báo cạnh. Luật C dùng cùng ngân
sách theo từng khách (RCE35).
% EV: sp/wf15/runner/engine.js:36-39 (B_r = round(alpha * slk)); kiểm toán 23/10 §7-W10, §7-W11; ADR-082 (chủ chốt EOHZ35 là đối chứng chính, 25/10/2026)

Bộ điều phối có hai kiểu. Kiểu chỉ chèn (INS) là kiểu của RideBound ở các mục trước. Kiểu gán lại (G2) cho phép, ở mỗi
quyết định, dời tối đa một khách đang chờ sang một trong 8 xe gần nhất, với giờ đón không muộn hơn giờ đón dự kiến của xe
cũ. Kiểu điều khiển theo điểm khả năng đi chung (Shadow) của bài cũng được cài lại: khách có điểm thấp nhận ngân sách 0,
khách còn lại giữ toàn bộ độ dư. Ngưỡng điểm là phân vị 72\% (INS) và 71\% (G2) của điểm trên các ngày khám phá, được hiệu
chỉnh để số khách phục vụ của Shadow với hàng rào nới gần bằng EO35; trên 60 ngày so sánh, thực tế 67,6--69,1\% khách được
phục vụ nhận ngân sách 0. Ba cách xử lý ngân sách đi với Shadow là SHF (hàng rào chặt), SHHZ (hàng rào nới) và SHRC (luật C).
% EV: sp/wf15/scripts/shlib.js:14-15 (g2Targets 8, ratchet); sp/wf15/SPEC.md:842, 1445; sp/wf15/work/CALIBRATION-RULE.md:24-34; sp/wf15/frozen/shadow-params.json
% EV (26/10): tỷ lệ khách bảo đảm thực tế: sp/wf15/results/CONFIRM-SHADOW-REPORT.txt:113-118 (S16 "realised guaranteed share" 67,64--69,09\% số khách được phục vụ); khóa ratchet của G2 so với giờ đón dự kiến theo kế hoạch cũ e, không phải giờ đã công bố: engine.js:19 (lock 'ratchet' <= e), CommitmentLockEvaluator.cs:183-186, 284-287

\subsection{Kết quả}
\label{muc:c4-via-ketqua}

Bảng~\ref{bang:c4-via} ghi giá trị trung bình mỗi ngày (tổng hai khung của trung bình ba lần rút mẫu) trên 60 ngày.

\begin{table}[htbp]
\centering
\caption[So sánh với Elmi Ordinary trên dữ liệu Via (bộ mô phỏng)]{Trung bình mỗi ngày trên 60 ngày, bộ mô phỏng.
Phục vụ: số khách được phục vụ. Đảo hứa: số khách có hai lần đổi giờ hiển thị liên tiếp ($\ge 60$ giây) ngược chiều nhau.
Dời: tỷ lệ khách được phục vụ bị điều phối lùi giờ trả muộn hơn ít nhất 5 giây. Trễ hạn: tỷ lệ khách được phục vụ đến muộn
hơn $p_0 + B_r + 60$ giây. Các cột phần trăm là trung bình của tỷ lệ theo ngày, nên lệch chút với tỷ lệ gộp nêu trong
lời văn.}
\label{bang:c4-via}
\footnotesize
\setlength{\tabcolsep}{3.2pt}
\begin{tabular}{@{}lrrrrrrrr@{}}
\toprule
& \multicolumn{4}{c}{Chỉ chèn (INS)} & \multicolumn{4}{c}{Gán lại (G2)}\\
\cmidrule(lr){2-5}\cmidrule(l){6-9}
Cách & Phục vụ & Đảo hứa & Dời (\%) & Trễ hạn (\%) & Phục vụ & Đảo hứa & Dời (\%) & Trễ hạn (\%)\\
\midrule
Không trần & 795,6 & 42,4 & 45,7 & -- & 804,5 & 88,2 & 47,1 & --\\
EO35 (chặt) & 723,2 & 21,6 & 31,4 & 1,12 & 730,6 & 39,6 & 32,5 & 1,03\\
EOHZ35 (nới) & 725,2 & 21,7 & 31,5 & 1,13 & 732,1 & 39,8 & 32,4 & 1,01\\
RCE35 (luật C) & 726,7 & 19,9 & 31,7 & 1,79 & 729,2 & 23,1 & 32,0 & 1,77\\
\midrule
SHF & 685,5 & 30,3 & 21,3 & 2,63 & 694,4 & 50,0 & 22,4 & 2,59\\
SHHZ & 723,9 & 31,4 & 20,6 & 3,27 & 732,8 & 53,6 & 21,9 & 3,17\\
SHRC & 719,8 & 30,9 & 20,5 & 3,25 & 723,3 & 36,9 & 21,4 & 3,10\\
\bottomrule
\end{tabular}
\end{table}
% EV: p5-via-adapter-2026-10-25/TABLES-2x3-60days.txt (tính lại độc lập; trùng các số ở kiểm toán 23/10 §6-P1 bước 4, §7-W1, W3, W4, W7)

So với hàng rào nới cùng ngân sách, khi chỉ chèn, luật C phục vụ ngang (+1,5 khách mỗi ngày, khoảng tin cậy 95\% theo
ngày từ +0,4 tới +2,6) và giảm đảo hứa ít (−1,7 mỗi ngày). Khi gán lại, luật C giảm 16,7 lần đảo hứa mỗi ngày (−42\%;
khoảng tin cậy −18,0 tới −15,3; 60 trên 60 ngày), đổi lại phục vụ ít hơn 2,8 khách mỗi ngày (0,4\%; khoảng tin cậy −4,4
tới −1,3). So với hàng rào chặt, hai hiệu tương ứng là −16,5 lần đảo hứa và −1,3 khách (khoảng tin cậy chứa 0). Ở mức độ
dư 35\% này, hai cách mở rộng hàng rào của Elmi cho kết quả gần như nhau, trừ số khách phục vụ (hàng rào nới nhiều hơn 1,5
tới 2,0 khách mỗi ngày).

Con số $-42\%$ cần đọc cùng hai mốc khác. Thứ nhất, chính việc gán lại sinh thêm đảo hứa: với hàng rào nới, gán lại nâng số
lần đảo hứa từ 21,7 lên 39,8 và số khách phục vụ từ 725,2 lên 732,1 mỗi ngày. Luật C khi gán lại có 23,1 lần đảo hứa và
phục vụ 729,2 khách, tức so với không gán lại vẫn thêm 1,4 lần đảo hứa nhưng phục vụ thêm 4,1 khách mỗi ngày: luật C bỏ
khoảng 92\% số lần đảo hứa do gán lại sinh ra và giữ khoảng 59\% phần phục vụ tăng thêm. Thứ hai, mức giảm đến từ chính cơ
chế của luật (tính độ lớn phần dời theo cả hai chiều và không hoàn lại), và đảo hứa được chọn làm chỉ số chính sau thử
nghiệm sơ bộ; kết quả cho biết luật làm đúng điều nó đặt ra, chưa cho biết khách được lợi. Độ lớn cũng đổi theo thiết lập:
trong các phép thăm dò khác, 16 ngày chưa từng chạy cho $-38{,}5\%$ (16 trên 16 ngày), thêm phí chuyển xe 300 giây cho
hàng rào còn $-28{,}6\%$, và ở mật độ 256 yêu cầu mỗi hai giờ mức giảm tính trên 1.000 khách chỉ còn khoảng một phần ba;
vì vậy nên đọc mức giảm là khoảng 28--43\%, không phải một con số cố định. Một giả thuyết rút ra sau khi xem số liệu, chưa
được kiểm, là riêng phí chuyển xe đã đạt khoảng 56\% mức giảm của luật C.
% EV: TABLES-2x3-60days.txt (INS EOHZ35 21,672 / 725,167; G2 EOHZ35 39,761 / 732,056; G2 RCE35 23,078 / 729,233 → +1,406 đảo, +4,066 khách; 92\% = 1 − 1,406/18,089; 59\% = 4,066/6,889); khớp wf17 b2-reassignment-practice (máy công ty)
% EV (máy công ty, chưa kiểm ở máy nhà): báo cáo wf17 RideBound-Do-Ben-Va-Gia-Tri-Loi-Ich-23-10.md §1, §2.3--2.5 (F16 −38,5\% 16/16; phí 300 s −28,6\%; P256 ≈ 1/3 tính trên 1.000 khách; 56\% là giả thuyết hậu kiểm R4)

Mức 35\% không phải điểm duy nhất cho hướng này: trong một phép thăm dò trên 10 ngày với $\alpha$ từ 10\% tới 100\%, khi gán
lại, luật C có ít hơn 33--55\% số lần đảo hứa so với hàng rào nới cùng $\alpha$, đổi lại phục vụ ít hơn 0,2--1,2\%; khi chỉ
chèn, hai cách gần như trùng nhau ở mọi $\alpha$. So ở cùng $\alpha$ chưa cho thấy hết cái giá. Khi chỉnh $\alpha$ của luật C
để nó phục vụ đúng bằng hàng rào nới (gán lại, $\alpha = 35\%$), luật C bớt 13,5 khách bị đảo hứa nhưng thêm 12,7 khách bị
điều phối dời giờ trả muộn quá một phút mỗi ngày; tính theo phút, tổng độ chênh giữa giờ đến và lời hứa đầu tăng khoảng 13,8
phút cho mỗi khách tránh được đảo hứa, trong đó 13,4 phút (khoảng 97\%) do điều phối. Như vậy luật C chỉ đáng dùng khi một lần bị đảo
hứa được coi là khó chịu ít nhất bằng $w^* = 0{,}94$ lần một lần bị dời muộn quá một phút (khoảng bootstrap theo ngày
0,46--1,61; $w^* = 0{,}77$ nếu đếm trễ theo mọi nguyên nhân); qua các mức $\alpha$ từ 10\% tới 75\%, ngưỡng này nằm trong
khoảng 0,3--1,2. Đây là một điều kiện, không phải kết quả về cảm nhận: chưa có số đo nào cho biết khách thật đặt hai loại
khó chịu này nặng nhẹ ra sao. Khi chỉ chèn, hai luật không phân biệt được. Một phép định giá minh họa dùng tham số giá trị
thời gian trong văn liệu cho thấy, với các tham số đó, luật C làm giảm 0,95--1,34\% doanh thu cước và chỉ hòa vốn nếu mỗi lần
đảo hứa tránh được đáng giá 12,8--17,8 phút ngồi xe; con số phụ thuộc giả định, không phải số đo.
% EV: p5-via-adapter-2026-10-25/N2-RESULT.md, N2-ALPHA-TABLE.tsv (10 ngày P5, chế độ plain, thăm dò)
% EV: breakeven-2026-10-26/BREAKEVEN-RESULT.md; BREAKEVEN-TABLE-v2.tsv:28 (dR −13,48; dLd +12,72; w* 0,944 [0,463; 1,614]; wstarLa 0,769; dDA_min 186,04 và dNet_min 180,28 → 13,8 và 13,4 phút mỗi khách, 97\%); BREAKEVEN-CURVES-v2.tsv (w* 0,316--1,241 qua α 10--75\%); kế hoạch phép hòa vốn viết sau một lần minh họa 0,77 trên cùng 10 ngày (docs/18 ADR-086); máy công ty chạy lại khớp (0,944)
% EV (máy công ty, chưa kiểm ở máy nhà): báo cáo wf17 §4.3--4.5 (định giá minh họa 0,95--1,34\% doanh thu cước; hòa vốn 12,8--17,8 phút ngồi xe mỗi lần đảo hứa tránh được)
% EV: TABLES-2x3-60days.txt (INS RCE35-EOHZ35 served 1.506 [0.369, 2.642]; visRev60 -1.728; G2 visRev60 -16.683 [-18.049, -15.317] 60/0; served -2.822 [-4.376, -1.269]; G2 RCE35-EO35 visRev60 -16.528, served -1.328 [-2.977, 0.321]; EOHZ35-EO35 served 2.006 và 1.494)

Mỗi luật giữ lời hứa theo định nghĩa của chính nó. Tính gộp trên 60 ngày, theo sổ của luật C, hàng rào chặt để tổng thay
đổi do điều phối vượt $B_r$ ở 1,18\% (chỉ chèn) và 4,17\% (gán lại) khách; theo hạn chót của hàng rào, luật C để 1,78\% và
1,75\% khách đến muộn quá $p_0 + B_r + 60$ giây, so với 1,11\% và 1,02\% của hàng rào chặt. Không luật nào hơn trên cả hai
thước đo.
% EV: tỷ lệ gộp trên 60 ngày (main.ndjson, confirm.ndjson); kiểm toán 23/10 §7-W1

Bảng~\ref{bang:c4-via-dinhnghia} tách hai loại số. Các số ở phần trên là hệ quả của chính định nghĩa luật: chúng chỉ cho
biết cài đặt làm đúng định nghĩa, không phải phát hiện. Các số ở phần dưới là kết quả đo, có thể khác đi khi đổi cấu hình.

\begin{table}[htbp]
\centering
\caption[Điều có theo định nghĩa và điều đo được (Via, bộ mô phỏng)]{Điều có theo định nghĩa của luật và điều đo được,
gộp trên 60 ngày (chỉ chèn / gán lại). Khách bảo đảm: khách có ngân sách 0 dưới Shadow.}
\label{bang:c4-via-dinhnghia}
\footnotesize
\begin{tabular}{@{}p{6.6cm}p{3.6cm}p{4.6cm}@{}}
\toprule
& Luật C & Hàng rào Elmi\\
\midrule
\multicolumn{3}{@{}l}{\textit{Theo định nghĩa (kiểm cài đặt)}}\\
Khách có tổng phần dời do điều phối vượt ngân sách & 0 (RCE35, SHRC) & 1,18\% / 4,17\% (EO35)\\
Khách bảo đảm bị điều phối đổi giờ trả ($\ge 1$ giây, hai chiều) & 0 (SHRC) & 0,56\% / 4,59\% (SHF); 0,65\% / 5,89\% (SHHZ)\\
Khách bảo đảm bị điều phối dời giờ trả muộn hơn ($\ge 1$ giây) & 0 (SHRC) & 0,56\% / 1,47\% (SHF); 0,65\% / 1,71\% (SHHZ)\\
\midrule
\multicolumn{3}{@{}l}{\textit{Đo được}}\\
Khách đến muộn hơn $p_0 + B_r + 60$ giây & 1,78\% / 1,75\% (RCE35) & 1,11\% / 1,02\% (EO35)\\
Khách bảo đảm có giờ hiển thị đổi $\ge 60$ giây & 4,86\% / 4,77\% (SHRC) & 4,20\% / 7,26\% (SHF)\\
Khách bảo đảm đến muộn hơn lời hứa đầu quá 60 giây & 4,34\% / 4,22\% (SHRC) & 3,55\% / 3,55\% (SHF)\\
Khách bảo đảm bị chuyển sang xe khác & 0 / 2 khách (SHRC) & 0 / 3.361 khách (SHF)\\
\bottomrule
\end{tabular}
\end{table}
% EV: tỷ lệ gộp trên 60 ngày từ confirm.ndjson (G_dispAny1, G_vis60, G_late60First, G_switched trên G_n) và main.ndjson (breachLedger, breachDeadline60 trên served); kiểm toán 23/10 §7-W5, W6
% EV (26/10): dòng "dời muộn hơn" = G_dispUp1 / G_n gộp 60 ngày từ confirm.ndjson: SHF 0,56\% / 1,47\%, SHHZ 0,65\% / 1,71\%, SHRC 0 (tính lại ở máy nhà; kiểm toán KT:639, 707)

Hai khách bảo đảm dưới SHRC bị chuyển sang xe khác dù tổng phần dời của họ bằng 0: luật C chỉ đếm việc dời giờ, không
cấm đổi xe. Kết quả bất lợi cho luật C cũng phải nêu đủ. Khi chỉ chèn, so với hàng rào nới, luật C để nhiều khách bị điều
phối lùi giờ trả hơn một chút (+0,24 điểm phần trăm, khoảng tin cậy +0,10 tới +0,39). Khi gán lại, nó phục vụ ít hơn có ý
nghĩa (−2,8 khách mỗi ngày). Theo hạn chót của hàng rào, nó để nhiều khách trễ hơn (bảng trên). Với Shadow, luật C phục vụ
ít hơn hàng rào nới 4,0 và 9,5 khách mỗi ngày, và khi gán lại có thêm 2,7 khách trễ quá 60 giây mỗi ngày (khoảng tin cậy
1,4 tới 4,0). Khách bảo đảm dưới luật C vẫn đến muộn quá 60 giây vì tắc đường, ở tỉ lệ cao hơn dưới hàng rào chặt. Khi chỉ
chèn, tổng độ lớn các lần giờ trả hiển thị thay đổi (bỏ qua thay đổi dưới 5 giây, cộng cho mọi khách) của luật C cao hơn hàng
rào chặt 12,5 phút mỗi ngày (khoảng tin cậy +4,2 tới +20,9); so với hàng rào nới thì hiệu không rõ (+5,6; khoảng tin cậy
$-2{,}6$ tới +13,9). Trong phần phân tích phụ của nghiên cứu Shadow, hai dự đoán ghi trước không khớp: khi gán lại, SHHZ
phục vụ nhiều hơn EO35 2,3 khách mỗi ngày (dự đoán gần 0), và SHRC có tổng độ lớn thay đổi hiển thị thấp hơn SHF 77 phút mỗi
ngày (dự đoán gần 0); còn khi chỉ chèn, SHRC phục vụ ít hơn EO35 3,3 khách mỗi ngày.
% EV: TABLES-2x3-60days.txt (INS RCE35-EOHZ35 pDispPct 0.240 [0.095, 0.385]; G2 served -2.822 [-4.376, -1.269]; SHRC71-SHHZ71 late60First 2.700 [1.363, 4.037]); confirm.ndjson G_switched SHRC71 = 2; kiểm toán 23/10 §6-P1 bước 7
% EV (26/10, tính lại ở máy nhà bằng tables_2x3.py trên main.ndjson + confirm.ndjson): INS RCE35-EO35 TV5min +12,504 [+4,150; +20,857] (21/39 ngày); INS RCE35-EOHZ35 TV5min +5,648 [−2,569; +13,865]; TV5 = metrics_sh.js:57-60 (ngưỡng 5 s); phân tích phụ wf15: CONFIRM-SHADOW-REPORT.txt:82 (S2 G2 SHHZ-EO35 served +2,283, pred ~0, sign p 0,022), :107 (S11 G2 SHRC-SHF TV5 −77,037, pred ~0, sign p 0,001), :55 (S4 INS SHRC-EO35 served −3,328, pred <0); kiểm toán KT:635

Với Shadow, luật C phục vụ ít hơn hàng rào nới 4,0 (chỉ chèn) và 9,5 (gán lại) khách mỗi ngày, đổi lại ít hơn 16,6 lần
đảo hứa mỗi ngày khi gán lại. Shadow với hàng rào chặt phục vụ ít hơn hẳn (SHF thấp hơn SHRC 29 tới 34 khách mỗi ngày).
Khoảng cách này đến từ cách hình thức hóa hàng rào chặt: khi giờ trả dự kiến $e$ của một khách có ngân sách 0 vượt $p_0$,
dù do làm tròn vị trí xe hay do tắc đường, xe của khách đó bị giữ nguyên tuyến và không nhận thêm khách. Nó không phải ưu
thế của luật C. Theo một phép thăm dò khác, khi giao thông phẳng, cho hàng rào chặt dung sai 1 giây xóa hết khoảng cách;
khi có tắc đường, khoảng cách còn lại một phần. Trên 10 ngày, dung sai 1 giây xóa 33\% (chỉ chèn) và 54\% (gán lại) khoảng
cách, 5 giây xóa 54\% và 66\%; trên 30 ngày của một mẫu khác, các tỉ lệ là 41\% và 65\% (1 giây), 58\% và 82\% (5 giây), và
theo tiêu chí ghi trước, trường hợp chỉ chèn chưa kết luận được. Với 60 giây, hàng rào chặt phục vụ nhiều hơn cả SHRC, đổi
lại điều phối dời giờ của 43 tới 66 khách bảo đảm mỗi ngày, so với 0 dưới SHRC. Các cách này nằm trên một đường đánh đổi
giữa số khách phục vụ và số khách bảo đảm bị dời; không cách nào hơn hẳn. So với cách siết đều dạng hàng rào nới, Shadow
giảm khoảng 10 điểm phần trăm số khách bị điều phối lùi giờ trả và có ít hơn hẳn số khách đến muộn hơn lời hứa đầu quá 60
giây ($-59$ khách mỗi ngày khi chỉ chèn, $-55$ khi gán lại, cả 60 trên 60 ngày), nhưng có nhiều đảo hứa hơn và nhiều khách
trễ hạn hơn.
% EV: TABLES-2x3-60days.txt (SHRC-SHHZ served -4.017, -9.544; visRev60 G2 -16.644; SHRC-SHF +34.350, +28.906; SHHZ-EOHZ35 pDispPct -10.86, -10.49; SHHZ-EOHZ35 late60First -59.033 (60/0), -55.378 (60/0)); kiểm toán 23/10 §2.4-F1, §7-W7
% EV: N1-RESULT.md (10 ngày: khoảng cách SHRC-SHF 35,0 → 23,6 → 16,1 (INS) và 29,4 → 13,6 → 10,0 (G2) ở 0 / 1 s / 5 s, tức 33\%/54\% và 54\%/66\%); (máy công ty, chưa kiểm ở máy nhà) báo cáo wf17 §2.2, r2: 30 ngày 41\%/65\% (1 s), 58\%/82\% (5 s), INS "chưa kết luận" theo tiêu chí ghi trước; giao thông phẳng 1 s xóa hết: kiểm toán KT:87-90

\subsection{Độ vững khi đổi cách định danh của bộ mô phỏng}
\label{muc:c4-via-dinhdanh}

Các kết quả trên được sinh khi bộ mô phỏng phá thế hòa giữa các phương án bằng số thứ tự đơn giản. RideBound thật phá thế
hòa bằng mã định danh băm của phương án. Để biết kết luận có phụ thuộc chi tiết này không, 60 kịch bản của 10 ngày (chọn
bằng hạt giống cố định trước khi chạy) được chạy lại với cách định danh của RideBound. Kết quả theo từng khách đổi ở hầu
hết kịch bản. Hiệu đảo hứa giữ dấu: với gán lại, luật C so với hàng rào nới giảm đảo hứa −14,9 (cách cũ) và −16,1 (cách
của RideBound) lần mỗi ngày; hiệu đảo hứa giữa Shadow với luật C và Shadow với hàng rào nới cũng giữ dấu. Các hiệu phục
vụ chỉ vài khách mỗi ngày thì không vững: trên mẫu 10 ngày này, luật C phục vụ ít hơn hàng rào nới 4,1 và 3,5 khách, nhưng
trên một mẫu 10 ngày khác, hiệu phục vụ đổi dấu khi đổi cách định danh (chỉ chèn từ −1,40 thành +2,10; gán lại từ −1,50
thành +1,20), nên được coi là nhiễu. Hai loại hiệu khác cũng không giữ dấu: số khách trễ quá 60 giây của luật C so với hàng
rào nới khi chỉ chèn, và các hiệu nhỏ giữa hàng rào nới và hàng rào chặt ngoài số khách phục vụ. Khóa luận không rút kết
luận nào từ các hiệu này.
% EV: p5-via-adapter-2026-10-25/P5-RESULT.md (cổng 0: 840/840 dòng trùng số liệu gốc; 7/84 ô đánh dấu, một ô chỉ do sai số dấu phẩy động)
% EV (máy công ty, chưa kiểm ở máy nhà): wf17/r1-identity-fence-pad/results/r1_tables.txt:5-6, 130-131 (mẫu 10 ngày khác: hiệu phục vụ INS −1,40 → +2,10, G2 −1,50 → +1,20)

\subsection{Kiểm chứng bộ mô phỏng bằng RideBound thật}
\label{muc:c4-cauno}

Bộ mô phỏng chỉ đáng tin tới mức nó làm đúng như RideBound. Vì vậy RideBound thật (bộ chạy C\# với FleetPy) và bộ mô phỏng
được chạy trên cùng các kịch bản Manhattan, và so từng khách: thời điểm hứa, giờ trả đã hứa, giờ trả thật, có được phục vụ,
tổng và số lần dời, đi chung. Dự đoán của bộ mô phỏng được khóa bằng mã băm trước khi chạy RideBound. Kết quả: 238 cặp
kịch bản--luật trùng hoàn toàn, gồm điều phối không trần, luật C trần cố định, trần từng lần, hạn chót chặt và nới (cả dạng
hạn chót M$^{+}$ với luật phục hồi ``không tệ hơn'' của chương này), ngân sách theo từng khách (luật C, hàng rào chặt và
nới), gán lại G2 có và không giới hạn số xe nhận, trên 21 kịch bản (11 kịch bản 108
yêu cầu, 8 xe; 10 kịch bản 128 yêu cầu, 10 xe). Để chạy được, lớp nối FleetPy được sửa ở hai chỗ, cả hai trong cùng một tình
huống: FleetPy đã khóa một chặng đón khách dài 0 giây mà RideBound gửi lại dưới dạng chưa bắt đầu. Bộ mô phỏng không đổi.
Phía RideBound, ngân sách theo từng khách và kiểu gán lại G2 được thêm vào trên một nhánh nghiên cứu, mặc định tắt, chỉ để
chạy phép kiểm này; chúng không thuộc phiên bản RideBound được đánh giá ở chương~\ref{chuong:phuongphap} và các mục trước
của chương này.
% EV: BAO-CAO-GIAI-DOAN-1.md (117/117, 11.521 khách); via-bridge-2026-10-22/BRIDGE-REPORT.md (55/55, 4.972 khách); ADR-081, ADR-082, ADR-084; bridge-phase1-2026-10-23/compare/nw-2026-10-25/final-compare.json (DLS-120 + no-worse 11/11, 976 khách); tổng 238 = 55 + 22 + 33 + 117 + 11

Trên các kịch bản Manhattan này, thời gian dừng là 0 giây (RideBound chưa mô hình xe bận phục vụ), và giới hạn 100 phương án
mỗi xe cùng ngân sách của bộ giải chưa lần nào ràng buộc (232.103 quyết định, tất cả tối ưu).
% EV: SOLVER-EVIDENCE-RESULT.md; ADR-083; kiểm toán 23/10 §7-W12

Ở cỡ Via, RideBound thật được chạy trên một họ kịch bản dựng từ dữ liệu Via ngày 13/2/2019, hai khung giờ, trong đó mỗi điểm
đón, trả được gộp về một điểm đại diện của khu vực TLC: 60 và 58 điểm, 509 và 508 yêu cầu, 62 xe. Phải gộp vì lần chạy
thật dựng đồ thị đầy đủ trên mọi điểm của kịch bản, và kịch bản gốc có khoảng 830 điểm. Ngân sách công việc của mỗi quyết
định được nâng lên, vì với 62 xe nó vượt mức mặc định và bộ chạy dừng có chủ đích; giới hạn cho từng xe giữ nguyên. Mỗi lần
chạy mất 1,4 tới 2 giờ, so với vài giây của bộ mô phỏng. Sáu cặp kịch bản--luật chạy xong và trùng bộ mô phỏng từng khách
(2.407 khách): hàng rào nới khi chỉ chèn và khi gán lại ở cả hai khung, luật C khi gán lại ở khung 17 giờ, và điều phối
không trần khi gán lại ở khung 8 giờ; số lần dời của G2 cũng trùng. Bốn lần chạy, gồm cả hai lần chạy luật C khi chỉ chèn,
đi hết mọi quyết định nhưng hết bộ nhớ khi ghi trạng thái cuối, nên không đọc được sổ của khách; hai lần chạy khác không được
thực hiện. Vì vậy, luật C khi chỉ chèn chưa được kiểm ở cỡ này; Shadow và dự báo vòng kín chỉ có trong bộ mô phỏng; và chính
các kịch bản 512 yêu cầu dùng ở mục~\ref{muc:c4-via-ketqua} chưa được chạy bằng RideBound. Các kết luận ở mục đó là kết luận
về bộ mô phỏng, đã được kiểm khớp với RideBound ở cỡ nhỏ và, cho các luật kể trên, ở cỡ Via trên kịch bản gộp theo khu.
% EV: RideBound-Code nhánh research/via-zone-2026-10-26, docs/18 ADR-085, ADR-086; kits/via-real-2026-10-26/VIA-LOG.md (V1–V6), compare/runA3-final-compare.json (6/6 cặp, 2.407/2.407 khách, số lần dời G2 4/4); máy công ty dựng lại từ TSV của kit, khớp

\subsection{Lịch sử thăm dò và giới hạn}
\label{muc:c4-via-gioihan}

Thiết lập Via (mật độ, số xe, thời gian đi tối đa, khung giờ) được chọn qua thăm dò để luật C có tình huống khác biệt.
Phép phân tích chính được viết và khóa bằng mã băm trước khi chạy, trong cùng hệ thống đã chạy thăm dò, không có dấu thời gian
bên ngoài; các giả thuyết được đặt sau thử nghiệm sơ bộ, không do người viết khóa luận đặt từ đầu; đảo hứa được chọn làm chỉ
số chính sau thử nghiệm sơ bộ; đối chứng chính của phần Shadow đã đổi từ SHHZ sang SHF rồi trở lại SHHZ sau thăm dò; nghiên
cứu Shadow dùng lại 60 ngày của nghiên cứu trước. Vì vậy mọi so sánh trong mục này là thăm dò. Mọi hiệu đo trong một cấu
hình (62 xe, 512 yêu cầu mỗi hai giờ, gấp 2 lần, $\alpha = 35\%$); ở cấu hình khác, hướng có thể giữ nhưng độ lớn đổi. Các
điểm đón, trả là điểm giả trong khu vực, tắc đường đều theo không gian, và thời gian dừng bằng 0.
% EV: kiểm toán 23/10 §4, §7-W19, §7-W20 (KT:389, 721), §6-N9 (không có đợt xác nhận mới: mọi so sánh trong mục ghi là thăm dò)

Các so sánh trong mục này không dùng dự báo. Trên cùng dữ liệu Via, một phép thăm dò khác cho thấy tác dụng của dự báo khác
nhau theo khung giờ: dùng dự báo ngay trong bộ điều phối làm kết quả trễ ở khung 17 giờ xấu đi (thêm khoảng 2,7 điểm phần
trăm so với cộng một khoảng đệm cùng độ dài vào lời hứa), và lời hứa không dự báo được hiệu chỉnh theo giờ từ các ngày trước
cho kết quả ngang với dự báo thô. Vì vậy mọi kết quả về dự báo cần được đọc theo từng khung giờ.
% EV (máy công ty, chưa kiểm ở máy nhà): báo cáo wf17 §0 ý 6, §5.1; wf17 r3 (dự báo vòng kín CL/CLF trên F14, F16; khung 17 giờ +2,7 điểm phần trăm so với nhánh đệm, khoảng +4 điểm so với lời hứa không dự báo)
% CHỜ CHỦ QUYẾT (kiểm toán §6-P1 bước 9, KT:637, 662): khai báo ai đặt giả thuyết và cách hình thức hóa đối chứng, theo quy định của khoa; chủ chốt với thầy rồi sửa câu "không do người viết khóa luận đặt từ đầu" cho nhất quán.
